% Editorial revision of the anonymous Law to Action manuscript.
\documentclass{article}
\PassOptionsToPackage{numbers,compress}{natbib}
\usepackage{neurips_2026_vericode}
\usepackage[utf8]{inputenc}
\usepackage[T1]{fontenc}
\usepackage{hyperref}
\usepackage{url}
\usepackage{booktabs}
\usepackage{amsfonts}
\usepackage{nicefrac}
\usepackage{microtype}
\usepackage{xcolor}
\usepackage{longtable,array}
\hypersetup{hidelinks,pdftitle={From Law to Action: Neuro-Symbolic Runtime Enforcement for MCP Agents},pdfauthor={}}
\title{From Law to Action: Neuro-Symbolic Runtime Enforcement for MCP Agents}
\workshoptitle{NeurIPS 2026 Workshop on AI for Verifiable Coding}
\author{Anonymous Author(s)}
\begin{document}
\maketitle

\begin{abstract}
Tool-using agents need to distinguish what a procedure recommends from what an action is authorized to do. We present a runtime enforcement architecture for Model Context Protocol (MCP) agents that separates source-derived constraints from delegated authority. Legal documents, vulnerability-and-fix records, and procedural skills produce distinct intermediate representations. Before a protected operation, the runtime checks the invocation context, scoped logical evidence, a signed capability, and durable consumption state. We evaluate this boundary in 900 fixed-action trials spanning 30 source families and five mechanism variants. On the held-out split, full enforcement produced no forbidden effects in 54 policy-forbidden attempts and completed all 54 allowed attempts. A policy-and-capability baseline allowed 18 of the 54 forbidden effects; unguarded execution allowed all 54. These results characterize enforcement of modeled policies at one sandbox handler. Legal interpretation accuracy and closed-loop model planning were not evaluated.
\end{abstract}

\section{Introduction}

An agent may select a useful procedure while overlooking a jurisdictional restriction, a known vulnerability, or the limits of its delegated authority. These concerns arise at different levels. A legal source describes obligations and exceptions; a vulnerability record identifies a potential security condition; a procedural skill describes intended work. None, by itself, determines whether a particular tool invocation should be permitted. The Model Context Protocol (MCP) provides a common interface for tool discovery and invocation, while applications remain responsible for enforcing their security requirements \citep{mcp,mcp-auth}.

We study an enforcement boundary immediately before a protected tool operation. The architecture compiles legal, security, and procedural sources into separate intermediate representations (IRs), binds their constraints to the exact proposed invocation, and checks the caller's capability before allowing an effect. The final decision depends on the current actor, arguments, resource, and environment. A previously acceptable request may therefore be denied after a change of audience, policy context, or use state.

For agentic coding, useful behavior must be considered alongside the effects a program or tool may produce. We isolate the enforcement boundary using fixed declarative actions and independently observed filesystem effects, measuring whether the mechanism prevents policy-forbidden mutations while preserving allowed work. This design separates enforcement from model planning and code generation.

We contribute an architecture separating source-derived constraints from delegated authority; a sandbox implementation combining capability verification, propositional checks, and durable consumption; and an evaluation of five matched mechanism variants across 30 source families.


\section{Related work}
\label{sec:related}

\paragraph{Runtime enforcement and agent security.}
Complete mediation requires every access to a protected resource to pass an authorization check \citep{saltzer1975protection}. Security automata, inlined monitors, and edit automata characterize policies that can be enforced during execution \citep{schneider2000enforceable,erlingsson2000irm,ligatti2005edit,hamlen2006enforcement}. Recent agent systems apply related ideas to tool use. AgentSpec supplies runtime triggers, predicates, and enforcement rules \citep{wang2025agentspec}; Progent checks symbolic tool privileges and policy updates \citep{shi2025progent}; and CaMeL uses capabilities to isolate agent data flow \citep{debenedetti2025camel}. Our design builds on this tradition by separating source-derived constraints from the capability needed for a concrete operation.

\paragraph{Executable law and checked evidence.}
Translating a statute into a formal representation entails an interpretation of the source \citep{sergot1986bna,prakken2015lawlogic,benzmueller2020logikey}. Catala makes legal specifications executable and has been applied to tax law \citep{catala,merigoux2021tax}. Legal IR serves a narrower role here: it represents selected facets of source documents for use in a runtime decision. Its semantic fidelity requires separate evaluation. Proof-carrying code and authentication place a checked proof before execution or authorization \citep{necula1997pcc,appel1999pca}, while translation validation and CompCert make the relationship between translation and verification explicit \citep{pnueli1998translation,leroy2009compcert}. We likewise distinguish the type of evidence a checker supplies. The evaluated propositional route establishes satisfiability results, not acceptance by a theorem-proving kernel.

\paragraph{Capabilities and policy languages.}
Object capabilities, attenuating caveats, relation-based authorization, and policy languages provide complementary ways to restrict authority \citep{miller2006robust,birgisson2014macaroons,pang2019zanzibar,cutler2024cedar}. UCAN specifies signed delegated capabilities \citep{ucan}. We use capability verification in both the lightweight baseline and the full enforcement mechanism; the latter additionally binds the decision to the current invocation and durable use state. This application-level check complements MCP's authorization protocol \citep{mcp-auth}.

\paragraph{Evaluation of coding and tool-using agents.}
Code-generation and issue-resolution benchmarks emphasize functional correctness and useful work \citep{chen2021codex,jimenez2024swebench}. Verified generation and autoformalization further connect programs or specifications to formal checks \citep{sun2024clover,yang2023leandojo,wu2022autoformalization}. AgentDojo, ToolEmu, InjecAgent, $\tau$-bench, and AgentHarm examine attacks, risky actions, policy following, and harmful tool use \citep{debenedetti2024agentdojo,ruan2024toolemu,zhan2024injecagent,yao2024taubench,andriushchenko2025agentharm}. These studies motivate measuring successful authorized work alongside forbidden effects. Our experiment compares mechanism variants within one implementation; the cited systems were not executed as external baselines.

\section{From source documents to action constraints}
\label{sec:sources}

The architecture uses three source streams because normative constraints, security evidence, and procedural intent answer different questions (Table~\ref{tab:1}). Each representation retains source identity, provenance, and declared limitations. Evidence from one stream does not confer authority on another.



\paragraph{Legal sources.}
A restricted grammar extracts seven fields: modality, actor, action, object, conditions, exceptions, and temporal qualifiers. Source maps connect the resulting representation to the document, and loss diagnostics identify unsupported constructs. This makes the compilation inspectable, but inspection alone does not establish whether the interpretation is legally correct. The evaluated legal cohort contains six documents; the broader source inventory is described in Appendix~\ref{app:corpus}.

\paragraph{Vulnerability and repair records.}
CVEfixes links vulnerability records to repositories and fixes \citep{cvefixes}. The adapter creates scoped Security IR declarations and rejects wildcard scope and inserted authority grants. Vulnerable/fixed pairs provide candidate contrasts for a restriction, but applying that restriction to a tool requires an explicit connection between the recorded behavior and the proposed effect. The evaluation retains 12 repository families and 24 cases. Their vulnerability polarity and transfer to the export handler remain unestablished; the fixed-action experiment instead tests an explicitly modeled policy.

\paragraph{Procedural skills.}
SkillCenter supplies procedural Markdown. A deterministic normalizer extracts goals, guards, effects, and control edges, retaining source spans and distinguishing grounded from inferred content. The reader opens source bundles read-only, disables extension loading, and treats embedded commands as data. The evaluated cohort contains 12 skill families, each with a base procedure and an adversarial variant. A statement that a procedure is ``permitted'' remains a source assertion; permission to execute comes from the capability check. No trained skill encoder is used in this evaluation.

\begin{table}[t]
\centering
\caption{The three source representations contribute constraints or intent to an action decision. Execution authority is checked separately.}
\label{tab:1}
\begin{tabular}{@{}>{\raggedright\arraybackslash}p{.18\linewidth} >{\raggedright\arraybackslash}p{.36\linewidth} >{\raggedright\arraybackslash}p{.40\linewidth}@{}}
\toprule
Representation & Source and retained information & Role in the decision \\
\midrule
Legal IR & Legal text: modality, parties, conditions, exceptions, and time & Represents the selected normative constraints and their applicability. \\[4pt]
Security IR & Vulnerability records and fixes: scope, effects, and candidate restrictions & Describes a potential prohibition and the evidence needed to connect it to an operation. \\[4pt]
Intent IR & Procedural skills: goals, actions, guards, effects, and control flow & Describes the proposed work and distinguishes source-grounded content from inference. \\
\bottomrule
\end{tabular}
\end{table}

Appendix~\ref{app:corpus} records source revisions, inspection counts, and sampling context for the 30 evaluated families.

\section{Context-bound runtime enforcement}
\label{sec:method}

The runtime connects source constraints to an operation through four stages: invocation binding, constraint checking, authorization, and guarded execution. Independent effect observation provides the measurement used in the experiment.

\paragraph{Invocation binding.}
An invocation envelope records the source identities, tenant, actor, audience, tool and version, arguments, proposed effects, environment, nonce, and relevant state roots. Trusted adapters connect the effect description to the operation being requested. All checks refer to this same envelope, so evidence for one invocation cannot be reused after changing its arguments or context without revalidation.

\paragraph{Constraint checking.}
The compiler combines the applicable Legal, Security, and Intent constraints. Every required constraint must pass; a favorable score elsewhere cannot compensate for a failed hard condition. The evaluated logical route encodes a Hoare-style security invariant in quantifier-free Boolean logic (QF\_BOOL). SymPy DPLL runs in a child process, and an independent exhaustive truth-table checker verifies the satisfiability result. This route supplies satisfiability evidence only. Requests requiring a kernel-checked theorem cannot be satisfied by that evidence. Appendix~\ref{app:logic} distinguishes this evaluated route from the broader set of registered representations.

\paragraph{Authorization and consumption.}
A candidate that satisfies the represented constraints must also carry a current grant for the actor, tool, resource, and bounds of the requested action. The implementation verifies signed UCAN delegation using Ed25519 and rechecks the live context immediately before delegation. Durable consumption records the permitted use in a file-backed DuckDB store accessed through a typed Quack owner. The mechanism rejects replay, expiry, a wrong audience, changed arguments, undeclared effects, and changes to the relevant roots. The relationship between persistence, identity, and network transport is detailed in Appendix~\ref{app:infrastructure}.

\paragraph{Protected operation and observation.}
The experiment enables enforcement explicitly around a sandbox JSON-export handler. The handler takes a declarative export instruction and performs a filesystem mutation. An independent filesystem journal records whether the effect occurred and whether allowed work succeeded. This observation is separate from the authorization decision: an \texttt{allow} response alone does not count as a successful or safe execution.

Appendix~\ref{app:trace} provides an implementation mapping and a worked trace.

\section{Evaluation}
\label{sec:eval}

On the held-out final split, full enforcement (A4) produced no policy-forbidden effects in $54$ attempts and completed all $54$ allowed attempts, with no false denials. Lightweight policy plus UCAN verification (A3) permitted $18/54$ forbidden effects; unguarded execution (A0) permitted $54/54$. Table~\ref{tab:3} presents these results alongside the full study population. The comparison measures enforcement of a frozen modeled policy at one sandbox handler.

\subsection{Study design and measurement}
\label{sec:eval-protocol}

The predeclared study comprises $30$ source-lineage families and $60$ fixed-action cases: six legal families, 12 CVE repository pairs, and 12 SkillCenter families \citep{cvefixes}. Each family contributes one policy-allowed and one policy-forbidden case. A lineage-disjoint split assigns $6/6/18$ families to development, calibration, and final evaluation. Five matched arms use identical candidate identifiers, arguments, and the same sandbox handler. A0 executes without a guard. A1 and A2 add inert prompt-text and retrieval-context labels as model-free equivalence controls. A3 combines lightweight policy with Ed25519 UCAN verification. A4 adds explicit \texttt{ENFORCE}, exact-context binding, a qualified QF\_BOOL SAT job, and file-backed DuckDB consumption.

Seeds $104729$, $104759$, and $104761$ yield $60\times5\times3=900$ scored cells. The independent unit is the lineage family; repeated seeds and the paired cases do not constitute additional independent families. Labels describe the frozen modeled policy and machine-checkable behavior contract. Effects are measured independently through filesystem-journal events from a handler that accepts declarative \texttt{export\_json} instructions. No scientific model calls or model-generated programs enter this study.

\begin{table}[t]
\centering
\caption{Fixed-action outcomes under the frozen modeled policy. Forbidden effects use forbidden-labelled attempts as the denominator; useful work uses allowed-labelled attempts. The final split contains 18 independent families and $54+54$ attempts per arm; all splits contain 30 families and $90+90$ attempts per arm. False denials are $0/54$ on the final split and $0/90$ overall for every arm. A1 and A2 are model-free controls.}
\label{tab:3}
\begin{tabular}{@{}lcccc@{}}
\toprule
& \multicolumn{2}{c}{Held-out final split} & \multicolumn{2}{c}{All splits} \\
\cmidrule(lr){2-3}\cmidrule(l){4-5}
Arm & \shortstack{Forbidden\\effects} & \shortstack{Allowed\\useful work} & \shortstack{Forbidden\\effects} & \shortstack{Allowed\\useful work} \\
\midrule
A0 unguarded & $54/54$ & $54/54$ & $90/90$ & $90/90$ \\[3pt]
A1 prompt label & $54/54$ & $54/54$ & $90/90$ & $90/90$ \\[3pt]
A2 retrieval label & $54/54$ & $54/54$ & $90/90$ & $90/90$ \\[3pt]
A3 policy + UCAN & $18/54$ & $54/54$ & $33/90$ & $90/90$ \\[3pt]
A4 \texttt{ENFORCE} & $0/54$ & $54/54$ & $0/90$ & $90/90$ \\
\bottomrule
\end{tabular}
\end{table}

\subsection{Effect prevention and useful work}
\label{sec:eval-admitted}

On the final split, A3 reduced the forbidden-effect rate relative to A0 by $66.7$ percentage points (paired family-bootstrap $95\%$ interval $[-88.9,-44.4]$). A4 reduced it by $100$ points relative to A0 ($[-100,-100]$) and by $33.3$ points relative to A3 ($[-55.6,-11.1]$). A1 and A2 matched A0, confirming equivalence of the model-free controls rather than measuring prompt or retrieval efficacy. The bootstrap resamples families within source types. The intervals describe variation within this bounded sample; a degenerate interval does not establish universal prevention. Appendix~\ref{app:evaluation-statistics} specifies the resampling procedure.

A3's $33$ forbidden effects across all splits involved wrong date or jurisdiction ($9$), forged receipts ($6$), replay ($6$), changed roots, clock, or environment ($6$), and undeclared handler effects ($6$). The final split contained six failures each for replay, changed roots/clock/environment, and wrong date/jurisdiction. A4 denied those same identities while completing all allowed attempts. This comparison identifies residual failures of the lightweight policy-and-capability arm; it does not isolate the contribution of each mechanism added by A4.

\subsection{Source-to-IR contracts}
\label{sec:eval-source}

All $60$ frozen cases retain a prediction, failure, or unsupported record, with no outcome-dependent cohort changes. Quoted source-span linkage held for $60/60$ cases, parse or schema validity for $50/60$, and scoped machine-contract coverage for $60/60$. Four CVE adapter errors and six policy-blocked skill rows remain in the denominators. An independent behavior-contract producer observed source differences for $24/24$ CVE cases, but polarity remained unknown for all $24$ cases ($12$ pairs), and transfer to the export handler was unsupported.

These are structural and behavioral contract checks. The skill-validator comparison ($18/18$) reuses the normalizer's Intent IR decoder and schema functions, so it measures shared-producer conformance rather than independent agreement. Expert legal applicability, exception fidelity, semantic accuracy, inter-annotator agreement, and legal validity were not measured; generated interpretations remain predictions. Table~\ref{tab:E1} summarizes the available evidence and its scope.

\subsection{Execution scope}
\label{sec:eval-negative}

All $900$ scored cells satisfied the study's resource qualification, with no scientific retries. Appendix~\ref{app:evaluation-details} documents startup recovery, resource measurements, and protocol amendments. Compact scored records contain effect counters but do not retain full per-cell proof or capability payloads.

Separate mediation, durable-consumption, and local-transport suites qualified the selected mechanisms. Development timings, mechanism ablations, and diagnostic replays are reported separately in the appendix and do not add scored observations. The planned closed-loop study was not executed and contributes no measured outcomes.

\section{Limitations and modeling assumptions}
\label{sec:limits}

The source representations are models of selected information. Legal IR omits unsupported constructs, and its interpretation has not been independently validated. Security IR expresses candidate prohibitions: a CVE/CWE identifier, fix, or source difference does not establish vulnerability absence or applicability to unrelated code. Skill Markdown remains untrusted input, including assertions of permission. Corpus release totals were not independently recounted and do not describe the evaluated population.

Mediation is confined to explicit \texttt{ENFORCE} around one sandbox export handler. Default \texttt{OFF} mode passes calls through; default in-memory consumption differs from the file-backed DuckDB configuration used to qualify restart behavior. A database commit alone does not establish exactly-once remote effects. Unprotected, observational, and nonidentical routes remain outside the safety claim, and local stdio and loopback HTTP/libp2p checks leave wide-area behavior untested.

The logical checking route supplies propositional satisfiability evidence, not a kernel-checked theorem. Omission of full per-cell proof and capability payloads limits retrospective inspection. The handler executes declarative instructions; closed-loop generation, trained normalizers, and learned encoders were not evaluated. The findings therefore concern the measured runtime boundary and automated contracts, without independent human semantic validation.


\section{Conclusion}

This study evaluates an action boundary that separates source-derived constraints from permission to execute. Across the fixed-action experiment, full context-bound enforcement prevented the tested policy-forbidden effects while preserving allowed work. The policy-and-capability baseline left residual failures involving context changes and replay, motivating checks at the point of execution. These observations support the evaluated configuration at one sandbox handler. Extending the result to legal interpretation or model-generated programs requires separate semantic and closed-loop evaluations.\label{la-main-text-end}

\clearpage
\bibliographystyle{unsrtnat}
\bibliography{related_work}
\clearpage
\appendix
\section{Corpus lineage and source provenance}
\label{app:corpus}
\setcounter{table}{0}
\renewcommand{\thetable}{\thesection.\arabic{table}}
\renewcommand{\theHtable}{\thesection.\arabic{table}}

The evaluated sample comprises six legal documents, 12 CVE repository families, and 12 SkillCenter families. Each family contributes two cases, giving 60 cases in total. The source inspection used to construct this sample read six legal documents, 4,329 rows from the first CVEfixes shard, and 1,182 SkillCenter index/content rows. These inspection counts are distinct from both the evaluated cases and the sizes reported on release cards. In particular, the release-card totals of 12,987 CVEfixes input rows and 216,972 SkillCenter physical records were not independently recounted. Municipal and national-legislation release totals likewise describe source collections, not evaluation coverage.

\begin{longtable}{@{}>{\raggedright\arraybackslash}p{.22\linewidth} >{\raggedright\arraybackslash}p{.34\linewidth} >{\raggedright\arraybackslash}p{.38\linewidth}@{}}
\caption{Legal source inventory and the provenance distinctions required for interpretation. The evaluated legal sample remains six documents.}\label{tab:A1}\\
\toprule
Source family & Inspected material & Provenance requirements \\
\midrule
\endfirsthead
\toprule
Source family & Inspected material & Provenance requirements \\
\midrule
\endhead
\bottomrule
\endlastfoot
Federal statutes & U.S.\ Code release with pinned revision & Provision identity, release point, amendments, and effective dates. \\[4pt]
State law & State statutes and a separate administrative-rule corpus & Jurisdiction, issuing authority, legal form, and temporal coverage. \\[4pt]
Federal regulatory material & Federal Register collection & Proposed or final status, publication date, effective date, and codification. \\[4pt]
Court rules and caselaw & State court rules and a Caselaw Access Project-derived collection & Court, jurisdiction, rule scope, and holdings versus quotations. \\[4pt]
Municipal and county law & Municipal harvests across releases & Publisher, jurisdiction, discovery and failure counts, and release aliases. \\[4pt]
International legislation & Netherlands, Belgian, and Argentine source collections & Language, authority, consolidation date, capture and extraction version, and coverage gaps. \\
\end{longtable}

\paragraph{Security-source revision.}
The inspected CVEfixes source is \texttt{hitoshura25/cvefixes} at revision
\begin{quote}\small\path{d4f5c4ea65329d9ccbb8a3b3149e5d06eda5edb2}.\end{quote}
The evaluated cohort contains 12 repository families and 24 cases. A further 84 controls were excluded and do not enter the empirical denominators. Source rows, graph nodes, edges, vectors, and unique vulnerability identifiers count different objects. Vulnerable/fixed source differences were observable in the retained records, but vulnerability polarity remained unknown under the behavior contract, and transfer to the export handler was not established.

\paragraph{Skill-source revision.}
The inspected SkillCenter source is \texttt{Tommysha/skillcenter-bundles} at revision
\begin{quote}\small\path{f9dd4fec3c86d85ebf116c7408ac5ce602c418a1}.\end{quote}
The evaluated skill cohort contains 12 families and 24 cases. SQLite bundles are read-only input containers; they are separate from the runtime's DuckDB state store. Index entries and available source records do not imply the same number of successfully normalized procedures.

\section{Logic representations and evidence types}
\label{app:logic}
\setcounter{table}{0}

The implementation separates a domain representation, its logical interpretation, and the backend that evaluates a supported encoding. The TDFOL and DCEC Python cores supply parsable native representations. Shared symbols or links between representations do not establish semantic equivalence. Checker outputs therefore retain their evidence type: a satisfiability result, a policy decision, a runtime observation, and a kernel-checked theorem are distinct outcomes. A timeout or an unavailable backend remains inconclusive.

\begin{longtable}{@{}>{\raggedright\arraybackslash}p{.23\linewidth} >{\raggedright\arraybackslash}p{.36\linewidth} >{\raggedright\arraybackslash}p{.35\linewidth}@{}}
\caption{Representation support and evaluation scope. The scored experiment used the QF\_BOOL route.}\label{tab:B1}\\
\toprule
Representation & Expressed information & Evaluation scope \\
\midrule
\endfirsthead
\toprule
Representation & Expressed information & Evaluation scope \\
\midrule
\endhead
\bottomrule
\endlastfoot
Propositional and first-order & Boolean operators, typed terms, predicates, and quantifiers & Support depends on the selected fragment and vocabulary. \\[4pt]
Frame and F-logic & Object identity, class membership, slots, and relational inheritance & Structured representation; no scored semantic-equivalence result. \\[4pt]
Deontic & Obligations, permissions, and prohibitions & Normative representation; execution permission still requires a capability. \\[4pt]
Temporal / TDFOL & Temporal operators, including always, eventually, until, and since & Representations registered; temporal backends were not scored. \\[4pt]
Event calculus & Events, fluents, and initiation or termination of state & Transition representation; not a scored backend. \\[4pt]
Datalog and policy rules & Bounded rules and authorization declarations & Rule derivation under an explicit open- or closed-world convention. \\[4pt]
Hoare-style / SAT & Preconditions, postconditions, and propositional invariants & QF\_BOOL checked by SymPy DPLL and an independent truth table; satisfiability evidence only. \\[4pt]
Interactive proof & Native proof terms and scripts & Kernel backends were unavailable in the evaluation environment. \\
\end{longtable}

QF\_LIA interpolation was not selected or available. SAT-only, simulated, policy, and monitor evidence cannot discharge a requirement for a kernel-checked theorem. Backend-specific availability and the distinction between scored runs and diagnostic checks are recorded in Appendix~\ref{app:evaluation-details}.

\section{State, artifacts, transport, and capabilities}
\label{app:infrastructure}
\setcounter{table}{0}

The implementation assigns separate responsibilities to mutable state, immutable artifacts, transport, and authorization (Table~\ref{tab:C1}). This separation determines what an observation at each layer can establish.

\begin{longtable}{@{}>{\raggedright\arraybackslash}p{.21\linewidth} >{\raggedright\arraybackslash}p{.37\linewidth} >{\raggedright\arraybackslash}p{.36\linewidth}@{}}
\caption{Infrastructure responsibilities and their boundaries.}\label{tab:C1}\\
\toprule
Layer & Responsibility & Boundary \\
\midrule
\endfirsthead
\toprule
Layer & Responsibility & Boundary \\
\midrule
\endhead
\bottomrule
\endlastfoot
Operational state & File-backed DuckDB, typed Quack commands, compare-and-swap revisions, and idempotency & Local committed state; remote exactly-once execution is not established. \\[4pt]
Source input & Read-only SQLite bundles and Parquet, JSON-LD, HTML, or WARC records & Storage format does not determine source meaning or execution authority. \\[4pt]
Retrieval & Sparse, vector, and graph products keyed by source identifiers & Retrieves candidate evidence; applicability requires a separate check. \\[4pt]
Immutable artifacts & Content-addressed blocks, manifests, maps, and decision records \citep{ipfs} & Identifies bytes; local identifiers do not demonstrate public availability. \\[4pt]
Peer transport & libp2p streams and framed MCP messages \citep{libp2p} & Carries requests; peer identity alone does not grant tool access. \\[4pt]
Delegated capability & Signed UCAN verification and revocation or replay state \citep{ucan} & Bounds delegated authority; legal interpretation requires separate validation. \\
\end{longtable}

A state client supplies a named command with typed parameters. After a lost reply, repeating an identical command can retrieve its committed result; reusing the same identity with a changed payload produces a conflict. The optional Quack gateway extension was unavailable in the isolated validation environment. The evaluated durable-consumption configuration uses the typed owner route described in Section~\ref{sec:method}.

UCAN checks cover issuer identity, delegation continuity, resource and ability bounds, temporal validity, attenuation, revocation, and replay policy. The separate transport checks exercise local processes and loopback routes. Wide-area libp2p behavior and an IPFS daemon were not evaluated. The auxiliary \texttt{AuthorizationGate} is a related policy gate rather than the wrapper evaluated for protected-effect prevention.

\section{Implementation mapping and worked trace}
\label{app:trace}
\setcounter{table}{0}

The supplement maps neutral aliases to included implementation definitions and records both original source commitments and hashes of the portable, path-normalized files. Private credentials and the live state store are excluded.

\begin{longtable}{@{}>{\raggedright\arraybackslash}p{.25\linewidth} >{\raggedright\arraybackslash}p{.69\linewidth}@{}}
\caption{Information inspected at each stage of the implemented route.}\label{tab:D1}\\
\toprule
Stage & Relevant inputs \\
\midrule
\endfirsthead
\toprule
Stage & Relevant inputs \\
\midrule
\endhead
\bottomrule
\endlastfoot
Legal compilation & Document revision, jurisdiction, source maps, selected facets, and unsupported constructs. \\[4pt]
Skill and CVE adaptation & Physical source records, procedural spans, vulnerable/fixed pairs, restriction scope, and uncertainty. \\[4pt]
Constraint checking & Candidate graph, applicable checks, evidence requirements, and current roots. \\[4pt]
Logical evaluation & QF\_BOOL jobs, provider results, independent truth-table results, and required evidence type. \\[4pt]
Guarded execution & Live invocation context, authenticated decision record, capability, expiry, and use state. \\[4pt]
Durable consumption & Command identity, generation, fence, revision, idempotency key, and committed result. \\
\end{longtable}

A separate diagnostic replay follows a permitted A4 skill case with seed 104729. The trace resolves source pins, derives Intent IR, binds the tool and arguments, and runs the QF\_BOOL check. The provider reports unsatisfiability and the independent checker agrees. After Ed25519 capability verification and durable consumption, the wrapper dispatches one JSON export; the independent journal records one event and useful work. The selected source adapter in this trace is the skill adapter, not the legal compiler.

Matched diagnostic replays with a wrong audience, an undeclared effect, or a reused invocation are denied at delegation or consumption and produce no forbidden effects. These traces explain the mechanism but are separate from the scored fixed-action executions. The supplement retains the original worked-trace identifier under the LA-020 record.

\section{Measurement coverage}
\label{app:coverage}
\setcounter{table}{0}

Table~\ref{tab:E1} distinguishes the primary runtime results from source checks, component qualification, and unmeasured questions. Qualification denotes a targeted implementation check rather than an additional experimental arm or an independent semantic evaluation.

\begin{longtable}{@{}>{\raggedright\arraybackslash}p{.26\linewidth} >{\raggedright\arraybackslash}p{.68\linewidth}@{}}
\caption{Evidence available for each measurement objective. Unmeasured studies are excluded from empirical claims.}\label{tab:E1}\\
\toprule
Objective & Evidence and scope \\
\midrule
\endfirsthead
\toprule
Objective & Evidence and scope \\
\midrule
\endhead
\bottomrule
\endlastfoot
Source-to-IR contracts & Across all 60 legal, CVE, and skill cases: source linkage $60/60$, parse/schema validity $50/60$, and machine-contract coverage $60/60$. Expert legal fidelity, applicability, exceptions, and agreement were not evaluated. \\[4pt]
CVE restriction specificity & 12 pairs and 24 cases, with 84 excluded controls. Source differences were observed in $24/24$ cases; polarity remained unknown in all 24. Transfer was not established. \\[4pt]
Skill grounding & Source-span linkage on all skill cases; $18/18$ shared-producer schema conformance. Six unsupported, policy-blocked rows remain in the accounting. \\[4pt]
Logical checks & QF\_BOOL SAT with an independent truth-table checker. Theorem kernels and QF\_LIA interpolation were unavailable. \\[4pt]
Protected effects & Full enforcement produced $0/90$ forbidden effects and $90/90$ allowed-task successes across the primary study. \\[4pt]
Durable-state behavior & $14/14$ targeted qualification cases covered generation/fence checks, concurrency, lost replies, and restart behavior. \\[4pt]
Transport parity & Local-process and loopback checks; wide-area behavior and an IPFS daemon were unavailable. \\[4pt]
Execution overhead & Development-only A4 instrumentation. Model-token counts and price-based comparisons were not measured. \\[4pt]
Closed-loop planning & The separate 900-cell planning study was not started: no model calls or generated programs. It contributes no outcome rates. \\
\end{longtable}

\section{Evaluation accounting and supplementary measurements}
\label{app:evaluation-details}

This appendix preserves the detailed accounting behind Section~\ref{sec:eval}. The primary scored population is the resource-qualified fixed-action matrix: $30$ independent source-lineage families, $60$ cases, five matched arms, and three seeds, for $900$ cells. The $6/6/18$ development/calibration/final split is by lineage. Each arm therefore has $90$ forbidden-labelled and $90$ allowed-labelled attempts overall ($180$ decisions), and $54$ of each on the final split. Model calls were $0$; no scientific generative model was pinned. Filesystem-journal observations measure the declarative handler's effects, not AST extraction from model-generated source. Independent human annotations were not collected.

\subsection{Statistical procedure}
\label{app:evaluation-statistics}
\setcounter{table}{0}

Table~\ref{tab:paired-effects} reports final-split paired differences in forbidden-effect rate. Values are expressed in percentage points and rounded to one decimal place; the retained analysis summary preserves machine precision. The stratified paired family bootstrap uses $2{,}000$ draws and seed $104729$, resampling families within the legal, CVE, and skill strata while retaining their paired cases and repeated seeds. The final split contains $3/7/8$ families in these strata, respectively. The independent unit remains the lineage family. Repeated seeds do not increase the number of independent families, and zero-width intervals do not establish absence of failures outside this population.

\begin{table}[ht]
\centering
\caption{Final-split paired effects and percentile bootstrap intervals (percentage points). Negative differences indicate fewer forbidden effects.}
\label{tab:paired-effects}
\begin{tabular}{@{}lrr@{}}
\toprule
Comparison & Difference & $95\%$ interval \\
\midrule
A1 versus A0 & $0.0$ & $[0.0,0.0]$ \\
A2 versus A0 & $0.0$ & $[0.0,0.0]$ \\
A3 versus A0 & $-66.7$ & $[-88.9,-44.4]$ \\
A4 versus A0 & $-100.0$ & $[-100.0,-100.0]$ \\
A4 versus A3 & $-33.3$ & $[-55.6,-11.1]$ \\
\bottomrule
\end{tabular}
\end{table}

\subsection{Execution provenance and resource accounting}
\label{app:evaluation-resources}

The primary summary is \path{results/fixed_actions_admitted/summary.json}, used with its matching cost report and claim guidance. The implementation commitment recorded for that matrix is \path{ee6d73c4a5fc30d05ec4e787fdeb46d6701a4f947b7c874aecabf603e31ce06c}. The compact rows retain independent effect counters; pinned executed mechanisms and a shared durable-store identity support the runtime observations. Full per-cell proof and capability payloads are not retained.

There were $902$ host attempts for $900$ scientific cells and $0$ scientific retries. Two original false host admissions, at schedule indices $459$ and $560$, remain explicitly labeled. Container-startup failures at indices $591$ and $625$ were recovered only after evidence established that scientific execution had not started; each recovered scientific cell was counted once. Measured descendant-inclusive group CPU time was $1608.817478$ seconds, below the frozen $18{,}000$-second stop. Peak cell memory was $274812928$ bytes, and the maximum individual-cell wall-to-exit time was $18.022665493073873$ seconds.

The startup recoveries used prospective cumulative-active allowances. Failed-startup plus recovered-run wall time totaled $11.766052$ seconds at index $591$ and $19.793533$ seconds at index $625$. Separate pauses were excluded from active runtime. These amendments do not establish compliance with an uninterrupted original absolute $20$-second interval.

\subsection{Supplementary studies and evidence mapping}
\label{app:evaluation-evidence-map}

The following identifiers map manuscript statements to the retained development records. They distinguish scored outcomes from qualification and diagnostic evidence.

\paragraph{Source contracts (LA-009).}
The $18/18$ skill-validator comparison uses the same Intent IR decoder and schema functions that the normalizer invokes. It is shared-producer conformance, not an independent check. The automated source-contract measures replace an uncollected expert legal-fidelity study. No independent legal applicability, exception-fidelity, semantic-accuracy, inter-annotator-agreement, or legal-validity measurements are available. 

\paragraph{Scored fixed-action matrix (LA-029).}
Only the matching admitted summary, cost report, and claim guidance supply the primary arm rates. The original compact fixed-action log and its rescue (LA-015) are protocol-unadmitted diagnostics. Historical tables and figures (LA-019) analyze that earlier log and cannot substitute for the admitted summary.

\paragraph{Ablations and comparisons (LA-016/LA-017).}
Paired mechanism ablations and their associated comparison wording are protocol-unadmitted diagnostic evidence. They do not establish scored model efficacy or additional primary arm rates.

\paragraph{Development efficiency study (LA-018).}
Cold, warm, stale-root, and stale-clock conditions each have $n=12$ A4 development observations. Mean end-to-end wall time was $0.012326$\,s for cold execution, $0.006635$\,s for warm execution, and $0.00322$\,s for each stale condition. Stale-root and stale-clock conditions recorded $12/12$ stale-evidence flags. Model tokens were unobserved, and no price conversion was made. These timings are instrumentation measurements, not a rescore of the $900$-cell matrix.

\paragraph{Execution traces (LA-020).}
Twelve-step traces replay frozen A4 identities for diagnosis. They retain intermediate artifacts separately from the compact scored rows and do not add scored cells.

\paragraph{Mechanism qualification.}
Mediation, DuckDB consumption, and local stdio/HTTP/libp2p parity suites passed their selected-claim cases. These are qualification results rather than arm-rate observations. The protected route wraps the export mutation in explicit \texttt{ENFORCE}. Its original implementation symbols are
\begin{quote}\raggedright\small
\path{SupervisorPreInvocationEnforcement.authorize_and_delegate}\\[2pt]
\path{BoundedExportHandler.export_json}.
\end{quote}
The supplement's \path{symbol_aliases.json} maps these definitions to neutral aliases.
Unprotected, observational, and nonidentical routes were exercised during qualification and excluded from the safety claim. Local transport qualification does not establish wide-area publication; an IPFS daemon was unavailable.

\paragraph{Proof authority and unavailable backends.}
The selected QF\_BOOL SAT route, checked independently by a truth table, carries satisfiability authority only. SAT/UNSAT cannot authorize an allow requiring \texttt{theorem\_proof}. Z3, cvc5, Vampire, E, Lean, Coq/Rocq, Isabelle, and QF\_LIA interpolation were absent from the sealed validation path and have no measured results in this study.

\paragraph{Closed-loop planning (protocol RQ3).}
The planned generated-code study was withdrawn without execution: $900/900$ scheduled cells were \texttt{not\_started}, model calls were $0$, and generated programs were $0$. These counts are scheduling accounting, not observed agent failures, and are not pooled with the fixed-action results. No trained skill normalizer or encoder was evaluated.

\section{Language-model assistance}
\label{app:disclosure}
Language-model assistance supported source inspection, implementation, execution orchestration, analysis of retained data, and manuscript preparation. Grok \texttt{grok-4.6} (xAI) supported drafting and inspection; Codex (GPT-6) supported implementation, recovery, analysis, and editorial revision. The fixed-action experiment made no scientific model calls: these assistants were not experimental arms. No independent human legal, security, or intent evaluation was collected.

\clearpage
\section*{NeurIPS Paper Checklist}

\begin{enumerate}

\item {\bf Claims}
    \item[] Question: Do the main claims made in the abstract and introduction accurately reflect the paper's contributions and scope?
    \item[] Answer: \answerYes{}
    \item[] Justification: The abstract and introduction describe the implemented enforcement boundary and the held-out results: A4 produced $0/54$ forbidden effects and completed $54/54$ allowed attempts. Legal interpretation accuracy and closed-loop planning remain outside the measured claims.
    \item[] Guidelines:
    \begin{itemize}
        \item The answer \answerNA{} means that the abstract and introduction do not include the claims made in the paper.
        \item The abstract and/or introduction should clearly state the claims made, including the contributions made in the paper and important assumptions and limitations. A \answerNo{} or \answerNA{} answer to this question will not be perceived well by the reviewers. 
        \item The claims made should match theoretical and experimental results, and reflect how much the results can be expected to generalize to other settings. 
        \item It is fine to include aspirational goals as motivation as long as it is clear that these goals are not attained by the paper. 
    \end{itemize}

\item {\bf Limitations}
    \item[] Question: Does the paper discuss the limitations of the work performed by the authors?
    \item[] Answer: \answerYes{}
    \item[] Justification: Section~\ref{sec:limits} discusses the source representations, the single protected handler, the scope of the propositional evidence, and the absence of human semantic validation or closed-loop evaluation.
    \item[] Guidelines:
    \begin{itemize}
        \item The answer \answerNA{} means that the paper has no limitation while the answer \answerNo{} means that the paper has limitations, but those are not discussed in the paper. 
        \item The authors are encouraged to create a separate ``Limitations'' section in their paper.
        \item The paper should point out any strong assumptions and how robust the results are to violations of these assumptions (e.g., independence assumptions, noiseless settings, model well-specification, asymptotic approximations only holding locally). The authors should reflect on how these assumptions might be violated in practice and what the implications would be.
        \item The authors should reflect on the scope of the claims made, e.g., if the approach was only tested on a few datasets or with a few runs. In general, empirical results often depend on implicit assumptions, which should be articulated.
        \item The authors should reflect on the factors that influence the performance of the approach. For example, a facial recognition algorithm may perform poorly when image resolution is low or images are taken in low lighting. Or a speech-to-text system might not be used reliably to provide closed captions for online lectures because it fails to handle technical jargon.
        \item The authors should discuss the computational efficiency of the proposed algorithms and how they scale with dataset size.
        \item If applicable, the authors should discuss possible limitations of their approach to address problems of privacy and fairness.
        \item While the authors might fear that complete honesty about limitations might be used by reviewers as grounds for rejection, a worse outcome might be that reviewers discover limitations that aren't acknowledged in the paper. The authors should use their best judgment and recognize that individual actions in favor of transparency play an important role in developing norms that preserve the integrity of the community. Reviewers will be specifically instructed to not penalize honesty concerning limitations.
    \end{itemize}

\item {\bf Theory assumptions and proofs}
    \item[] Question: For each theoretical result, does the paper provide the full set of assumptions and a complete (and correct) proof?
    \item[] Answer: \answerNA{}
    \item[] Justification: The paper presents an empirical systems study and makes no theoretical claims requiring a proof. The evaluated QF\_BOOL route supplies satisfiability evidence.
    \item[] Guidelines:
    \begin{itemize}
        \item The answer \answerNA{} means that the paper does not include theoretical results. 
        \item All the theorems, formulas, and proofs in the paper should be numbered and cross-referenced.
        \item All assumptions should be clearly stated or referenced in the statement of any theorems.
        \item The proofs can either appear in the main paper or the supplemental material, but if they appear in the supplemental material, the authors are encouraged to provide a short proof sketch to provide intuition. 
        \item Inversely, any informal proof provided in the core of the paper should be complemented by formal proofs provided in appendix or supplemental material.
        \item Theorems and Lemmas that the proof relies upon should be properly referenced. 
    \end{itemize}

    \item {\bf Experimental result reproducibility}
    \item[] Question: Does the paper fully disclose all the information needed to reproduce the main experimental results of the paper to the extent that it affects the main claims and/or conclusions of the paper (regardless of whether the code and data are provided or not)?
    \item[] Answer: \answerYes{}
    \item[] Justification: Section~\ref{sec:eval} and Appendix~\ref{app:evaluation-details} specify the protocol, arms, seeds, and implementation commitment. The anonymous supplement records the environment and provides a command to reproduce the analysis from hashed retained observations; this does not constitute a fresh experimental run.
    \item[] Guidelines:
    \begin{itemize}
        \item The answer \answerNA{} means that the paper does not include experiments.
        \item If the paper includes experiments, a \answerNo{} answer to this question will not be perceived well by the reviewers: Making the paper reproducible is important, regardless of whether the code and data are provided or not.
        \item If the contribution is a dataset and\slash or model, the authors should describe the steps taken to make their results reproducible or verifiable. 
        \item Depending on the contribution, reproducibility can be accomplished in various ways. For example, if the contribution is a novel architecture, describing the architecture fully might suffice, or if the contribution is a specific model and empirical evaluation, it may be necessary to either make it possible for others to replicate the model with the same dataset, or provide access to the model. In general. releasing code and data is often one good way to accomplish this, but reproducibility can also be provided via detailed instructions for how to replicate the results, access to a hosted model (e.g., in the case of a large language model), releasing of a model checkpoint, or other means that are appropriate to the research performed.
        \item While NeurIPS does not require releasing code, the conference does require all submissions to provide some reasonable avenue for reproducibility, which may depend on the nature of the contribution. For example
        \begin{enumerate}
            \item If the contribution is primarily a new algorithm, the paper should make it clear how to reproduce that algorithm.
            \item If the contribution is primarily a new model architecture, the paper should describe the architecture clearly and fully.
            \item If the contribution is a new model (e.g., a large language model), then there should either be a way to access this model for reproducing the results or a way to reproduce the model (e.g., with an open-source dataset or instructions for how to construct the dataset).
            \item We recognize that reproducibility may be tricky in some cases, in which case authors are welcome to describe the particular way they provide for reproducibility. In the case of closed-source models, it may be that access to the model is limited in some way (e.g., to registered users), but it should be possible for other researchers to have some path to reproducing or verifying the results.
        \end{enumerate}
    \end{itemize}


\item {\bf Open access to data and code}
    \item[] Question: Does the paper provide open access to the data and code, with sufficient instructions to faithfully reproduce the main experimental results, as described in supplemental material?
    \item[] Answer: \answerYes{}
    \item[] Justification: The anonymous supplement includes all $9{,}933$ retained reducer inputs, the executed controller, driver, cell, and reducer sources, and an offline analysis command. Original and portable source hashes are distinguished. Native binaries, private keys, and source corpus bodies are excluded, so the supplied route reproduces the retained-data analysis.
    \item[] Guidelines:
    \begin{itemize}
        \item The answer \answerNA{} means that paper does not include experiments requiring code.
        \item Please see the NeurIPS code and data submission guidelines (\url{https://neurips.cc/public/guides/CodeSubmissionPolicy}) for more details.
        \item While we encourage the release of code and data, we understand that this might not be possible, so \answerNo{} is an acceptable answer. Papers cannot be rejected simply for not including code, unless this is central to the contribution (e.g., for a new open-source benchmark).
        \item The instructions should contain the exact command and environment needed to run to reproduce the results. See the NeurIPS code and data submission guidelines (\url{https://neurips.cc/public/guides/CodeSubmissionPolicy}) for more details.
        \item The authors should provide instructions on data access and preparation, including how to access the raw data, preprocessed data, intermediate data, and generated data, etc.
        \item The authors should provide scripts to reproduce all experimental results for the new proposed method and baselines. If only a subset of experiments are reproducible, they should state which ones are omitted from the script and why.
        \item At submission time, to preserve anonymity, the authors should release anonymized versions (if applicable).
        \item Providing as much information as possible in supplemental material (appended to the paper) is recommended, but including URLs to data and code is permitted.
    \end{itemize}


\item {\bf Experimental setting/details}
    \item[] Question: Does the paper specify all the training and test details (e.g., data splits, hyperparameters, how they were chosen, type of optimizer) necessary to understand the results?
    \item[] Answer: \answerYes{}
    \item[] Justification: Section~\ref{sec:eval-protocol} specifies the lineage splits, seeds, matched arms, frozen modeled policy, protected handler, and independent effect measurements. Appendix~\ref{app:evaluation-resources} and the supplement document the execution environment and resource accounting.
    \item[] Guidelines:
    \begin{itemize}
        \item The answer \answerNA{} means that the paper does not include experiments.
        \item The experimental setting should be presented in the core of the paper to a level of detail that is necessary to appreciate the results and make sense of them.
        \item The full details can be provided either with the code, in appendix, or as supplemental material.
    \end{itemize}

\item {\bf Experiment statistical significance}
    \item[] Question: Does the paper report error bars suitably and correctly defined or other appropriate information about the statistical significance of the experiments?
    \item[] Answer: \answerYes{}
    \item[] Justification: Section~\ref{sec:eval-admitted} reports paired family-level effect differences and bootstrap intervals. Appendix~\ref{app:evaluation-statistics} specifies the stratification, $2{,}000$ draws, seed, and treatment of repeated cases. Degenerate intervals are interpreted within the observed sample.
    \item[] Guidelines:
    \begin{itemize}
        \item The answer \answerNA{} means that the paper does not include experiments.
        \item The authors should answer \answerYes{} if the results are accompanied by error bars, confidence intervals, or statistical significance tests, at least for the experiments that support the main claims of the paper.
        \item The factors of variability that the error bars are capturing should be clearly stated (for example, train/test split, initialization, random drawing of some parameter, or overall run with given experimental conditions).
        \item The method for calculating the error bars should be explained (closed form formula, call to a library function, bootstrap, etc.)
        \item The assumptions made should be given (e.g., Normally distributed errors).
        \item It should be clear whether the error bar is the standard deviation or the standard error of the mean.
        \item It is OK to report 1-sigma error bars, but one should state it. The authors should preferably report a 2-sigma error bar than state that they have a 96\% CI, if the hypothesis of Normality of errors is not verified.
        \item For asymmetric distributions, the authors should be careful not to show in tables or figures symmetric error bars that would yield results that are out of range (e.g., negative error rates).
        \item If error bars are reported in tables or plots, the authors should explain in the text how they were calculated and reference the corresponding figures or tables in the text.
    \end{itemize}

\item {\bf Experiments compute resources}
    \item[] Question: For each experiment, does the paper provide sufficient information on the computer resources (type of compute workers, memory, time of execution) needed to reproduce the experiments?
    \item[] Answer: \answerYes{}
    \item[] Justification: Appendix~\ref{app:evaluation-resources} reports descendant-inclusive CPU time, peak cell memory, host-attempt counts, and individual-cell wall time. Startup recovery and excluded pause intervals are documented separately.
    \item[] Guidelines:
    \begin{itemize}
        \item The answer \answerNA{} means that the paper does not include experiments.
        \item The paper should indicate the type of compute workers CPU or GPU, internal cluster, or cloud provider, including relevant memory and storage.
        \item The paper should provide the amount of compute required for each of the individual experimental runs as well as estimate the total compute. 
        \item The paper should disclose whether the full research project required more compute than the experiments reported in the paper (e.g., preliminary or failed experiments that didn't make it into the paper). 
    \end{itemize}
    
\item {\bf Code of ethics}
    \item[] Question: Does the research conducted in the paper conform, in every respect, with the NeurIPS Code of Ethics \url{https://neurips.cc/public/EthicsGuidelines}?
    \item[] Answer: \answerYes{}
    \item[] Justification: The study uses public legal texts and retrieval-only vulnerability and skill sources, involves no human subjects, preserves submission anonymity, and does not redistribute source bodies or exploit payloads.
    \item[] Guidelines:
    \begin{itemize}
        \item The answer \answerNA{} means that the authors have not reviewed the NeurIPS Code of Ethics.
        \item If the authors answer \answerNo, they should explain the special circumstances that require a deviation from the Code of Ethics.
        \item The authors should make sure to preserve anonymity (e.g., if there is a special consideration due to laws or regulations in their jurisdiction).
    \end{itemize}


\item {\bf Broader impacts}
    \item[] Question: Does the paper discuss both potential positive societal impacts and negative societal impacts of the work performed?
    \item[] Answer: \answerNo{}
    \item[] Justification: The paper discusses the technical limits of runtime enforcement but does not provide a dedicated analysis of positive and negative societal impacts.
    \item[] Guidelines:
    \begin{itemize}
        \item The answer \answerNA{} means that there is no societal impact of the work performed.
        \item If the authors answer \answerNA{} or \answerNo, they should explain why their work has no societal impact or why the paper does not address societal impact.
        \item Examples of negative societal impacts include potential malicious or unintended uses (e.g., disinformation, generating fake profiles, surveillance), fairness considerations (e.g., deployment of technologies that could make decisions that unfairly impact specific groups), privacy considerations, and security considerations.
        \item The conference expects that many papers will be foundational research and not tied to particular applications, let alone deployments. However, if there is a direct path to any negative applications, the authors should point it out. For example, it is legitimate to point out that an improvement in the quality of generative models could be used to generate Deepfakes for disinformation. On the other hand, it is not needed to point out that a generic algorithm for optimizing neural networks could enable people to train models that generate Deepfakes faster.
        \item The authors should consider possible harms that could arise when the technology is being used as intended and functioning correctly, harms that could arise when the technology is being used as intended but gives incorrect results, and harms following from (intentional or unintentional) misuse of the technology.
        \item If there are negative societal impacts, the authors could also discuss possible mitigation strategies (e.g., gated release of models, providing defenses in addition to attacks, mechanisms for monitoring misuse, mechanisms to monitor how a system learns from feedback over time, improving the efficiency and accessibility of ML).
    \end{itemize}
    
\item {\bf Safeguards}
    \item[] Question: Does the paper describe safeguards that have been put in place for responsible release of data or models that have a high risk for misuse (e.g., pre-trained language models, image generators, or scraped datasets)?
    \item[] Answer: \answerNA{}
    \item[] Justification: The submission releases no pretrained generative model or scraped web-scale dataset. Source bodies remain retrieval-only, with provenance records and hashes.
    \item[] Guidelines:
    \begin{itemize}
        \item The answer \answerNA{} means that the paper poses no such risks.
        \item Released models that have a high risk for misuse or dual-use should be released with necessary safeguards to allow for controlled use of the model, for example by requiring that users adhere to usage guidelines or restrictions to access the model or implementing safety filters. 
        \item Datasets that have been scraped from the Internet could pose safety risks. The authors should describe how they avoided releasing unsafe images.
        \item We recognize that providing effective safeguards is challenging, and many papers do not require this, but we encourage authors to take this into account and make a best faith effort.
    \end{itemize}

\item {\bf Licenses for existing assets}
    \item[] Question: Are the creators or original owners of assets (e.g., code, data, models), used in the paper, properly credited and are the license and terms of use explicitly mentioned and properly respected?
    \item[] Answer: \answerYes{}
    \item[] Justification: The paper cites the source collections and records their revisions. The supplement documents GovInfo U.S.\ government works, Apache-2.0 CVEfixes packaging, and MIT SkillCenter packaging. Upstream code and nested sources retain their own licensing requirements.
    \item[] Guidelines:
    \begin{itemize}
        \item The answer \answerNA{} means that the paper does not use existing assets.
        \item The authors should cite the original paper that produced the code package or dataset.
        \item The authors should state which version of the asset is used and, if possible, include a URL.
        \item The name of the license (e.g., CC-BY 4.0) should be included for each asset.
        \item For scraped data from a particular source (e.g., website), the copyright and terms of service of that source should be provided.
        \item If assets are released, the license, copyright information, and terms of use in the package should be provided. For popular datasets, \url{paperswithcode.com/datasets} has curated licenses for some datasets. Their licensing guide can help determine the license of a dataset.
        \item For existing datasets that are re-packaged, both the original license and the license of the derived asset (if it has changed) should be provided.
        \item If this information is not available online, the authors are encouraged to reach out to the asset's creators.
    \end{itemize}

\item {\bf New assets}
    \item[] Question: Are new assets introduced in the paper well documented and is the documentation provided alongside the assets?
    \item[] Answer: \answerYes{}
    \item[] Justification: The anonymous supplement documents retained observations, source and runtime provenance, analysis instructions, measurement-integrity controls, and source-retrieval procedures. It explains the limits of the included artifacts and excludes private keys and the live state store.
    \item[] Guidelines:
    \begin{itemize}
        \item The answer \answerNA{} means that the paper does not release new assets.
        \item Researchers should communicate the details of the dataset\slash code\slash model as part of their submissions via structured templates. This includes details about training, license, limitations, etc. 
        \item The paper should discuss whether and how consent was obtained from people whose asset is used.
        \item At submission time, remember to anonymize your assets (if applicable). You can either create an anonymized URL or include an anonymized zip file.
    \end{itemize}

\item {\bf Crowdsourcing and research with human subjects}
    \item[] Question: For crowdsourcing experiments and research with human subjects, does the paper include the full text of instructions given to participants and screenshots, if applicable, as well as details about compensation (if any)? 
    \item[] Answer: \answerNA{}
    \item[] Justification: No crowdsourcing or human-subjects study was conducted. Human semantic and legal-fidelity measurements are not reported.
    \item[] Guidelines:
    \begin{itemize}
        \item The answer \answerNA{} means that the paper does not involve crowdsourcing nor research with human subjects.
        \item Including this information in the supplemental material is fine, but if the main contribution of the paper involves human subjects, then as much detail as possible should be included in the main paper. 
        \item According to the NeurIPS Code of Ethics, workers involved in data collection, curation, or other labor should be paid at least the minimum wage in the country of the data collector. 
    \end{itemize}

\item {\bf Institutional review board (IRB) approvals or equivalent for research with human subjects}
    \item[] Question: Does the paper describe potential risks incurred by study participants, whether such risks were disclosed to the subjects, and whether Institutional Review Board (IRB) approvals (or an equivalent approval/review based on the requirements of your country or institution) were obtained?
    \item[] Answer: \answerNA{}
    \item[] Justification: No human-subjects study was conducted, and no IRB or equivalent approval was obtained.
    \item[] Guidelines:
    \begin{itemize}
        \item The answer \answerNA{} means that the paper does not involve crowdsourcing nor research with human subjects.
        \item Depending on the country in which research is conducted, IRB approval (or equivalent) may be required for any human subjects research. If you obtained IRB approval, you should clearly state this in the paper. 
        \item We recognize that the procedures for this may vary significantly between institutions and locations, and we expect authors to adhere to the NeurIPS Code of Ethics and the guidelines for their institution. 
        \item For initial submissions, do not include any information that would break anonymity (if applicable), such as the institution conducting the review.
    \end{itemize}

\item {\bf Declaration of LLM usage}
    \item[] Question: Does the paper describe the usage of LLMs if it is an important, original, or non-standard component of the core methods in this research? Note that if the LLM is used only for writing, editing, or formatting purposes and does \emph{not} impact the core methodology, scientific rigor, or originality of the research, declaration is not required.
    %this research? 
    \item[] Answer: \answerNA{}
    \item[] Justification: LLMs were not part of the scored mechanism, and the fixed-action experiment made no scientific model calls. Appendix~\ref{app:disclosure} discloses Grok and Codex assistance with implementation, orchestration, analysis, and writing.
    \item[] Guidelines:
    \begin{itemize}
        \item The answer \answerNA{} means that the core method development in this research does not involve LLMs as any important, original, or non-standard components.
        \item Please refer to our LLM policy in the NeurIPS handbook for what should or should not be described.
    \end{itemize}

\end{enumerate}

\end{document}
